\documentclass{article}
\usepackage[utf8]{inputenc}
\usepackage{graphicx}
\usepackage{booktabs}
\usepackage{hyperref}

\title{Function-Morphology Correspondence in Neural Networks: Constrained Geometry Induces Causal Motif Specialization}
\author{Research Team}
\date{\today}

\begin{document}
\maketitle

\begin{abstract}
Recent advances in neural architectures often rely on uniform scaling of homogeneous transformer blocks. In this paper, we investigate the Finite Operator Grammar (FOG) hypothesis... (Refer to markdown draft for full abstract). We provide strong evidence that FOG mitigates capacity starvation and catastrophic interference, demonstrating sharply isolated functional pathways via causal ablation.
\end{abstract}

\section{Introduction}
The uniform transformer architecture relies on polysemanticity to multiplex distinct computational operations within shared dimensional subspaces...

\section{Hypothesis and Framework}
We hypothesize a Function-Morphology Correspondence: specific cognitive motifs demand specialized subspace geometries.

\section{Experimental Setup}
We compare our heterogeneous FOG architecture against Parameter-Matched Uniform Baselines across multiple regimes...

\section{Results}
\subsection{Causal Specialization}
Crucially, causal knockout analysis reveals that FOG naturally isolates operations. Ablating a specific head in FOG catastrophically degrades \texttt{Copy} performance while leaving \texttt{Retrieval} virtually untouched.

\section{Limitations}
Our findings are constrained to small-scale algorithmic tasks and theory-guided interpretations of motif signatures.

\section{Conclusion}
We provide strong evidence that enforcing structural heterogeneity mitigates catastrophic interference.

\end{document}
